\chapter{Falsifiable Predictions}
\label{ch:predictions}

SDT is a theory that makes specific, falsifiable predictions. If any of the following predictions are shown to be incorrect, the theory requires revision.

\section{Atomic and Nuclear Predictions}

\begin{enumerate}
  \item \textbf{Koppa universality:} No element with $Z > 100$ should yield $\kop \neq 0.5464$ when relativistic corrections are properly applied.
  \item \textbf{Screening saturation:} The per-electron screening efficiency $\sigma/(N{-}1)$ must plateau near $0.93$ for all heavy elements with filled d and f shells.
  \item \textbf{Screening from geometry:} The screening function $\sigma(Z,N)$ should be derivable from solid-angle geometry alone, without empirical fitting.
  \item \textbf{d-shell transition:} The $12\%$ jump in screening efficiency at $N \approx 28$ should correlate with X-ray scattering cross-section changes at the same electron count.
  \item \textbf{Neutron structure:} The neutron magnetic moment should be derivable from the proton-electron composite model with geometric corrections.
\end{enumerate}

\section{Gravitational Predictions}

\begin{enumerate}
  \item \textbf{Polar radius principle:} Using the polar radius of any oblate body should improve orbital velocity predictions vs.\ mean or equatorial radius.
  \item \textbf{Stellar rotation formula:} $\kop^2 = \pi(c/v_{\text{rot}})$ should hold for main-sequence stars but fail for planets, white dwarfs, and neutron stars.
  \item \textbf{No graviton:} Gravitational waves are pressure waves in the spation medium. No graviton particle will ever be detected.
\end{enumerate}

\section{Cosmological Predictions}

\begin{enumerate}
  \item \textbf{No dark matter detection:} All direct, indirect, and collider searches for dark matter particles will continue to return null results.
  \item \textbf{No dark energy:} Future measurements will confirm that the ``accelerating expansion'' is consistent with geometric energy dilution without a cosmological constant.
  \item \textbf{CMB as horizon:} The CMB temperature should show a distance-dependent profile consistent with diluted galactic emission, not a single-temperature blackbody.
\end{enumerate}

\section{Materials Science Predictions (Atomicus)}

\begin{enumerate}
  \item \textbf{Crystal structure from $\kop$:} The preferred crystal structure of any element should be predictable from its valence screening parameters.
  \item \textbf{Band gap from screening:} Semiconductor band gaps should correlate with the screening regime transition for the valence electrons.
  \item \textbf{Superconductivity:} The critical temperature $T_c$ should correlate with the geometric stability of the d-shell screenin configuration.
\end{enumerate}
